% ============================================================
%  MTH636 -- Game Theory | Lecture 3
%  Dominated Strategies & Iterated Elimination
% ============================================================
\documentclass[12pt, a4paper]{article}

% ============================================================
%  MTH636 -- Game Theory
%  Notes Preamble — Color-Coded, TikZ-Enhanced
% ============================================================

% --- Core packages ---
\usepackage[T1]{fontenc}
\usepackage[utf8]{inputenc}
\usepackage{lmodern}
\usepackage[margin=2.5cm, top=2.8cm, bottom=2.8cm]{geometry}
\usepackage{microtype}

% --- Math ---
\usepackage{amsmath, amssymb, amsthm, mathtools}
\usepackage{bm}        % bold math
\usepackage{cancel}    % strike-through in math

% --- Color & Boxes ---
\usepackage[dvipsnames, svgnames, x11names, table]{xcolor}
\usepackage[most]{tcolorbox}
\tcbuselibrary{breakable, skins, theorems}

% --- TikZ & Diagrams ---
\usepackage{tikz}
\usepackage{tikz-qtree}
\usetikzlibrary{arrows.meta, positioning, shapes.geometric, fit, backgrounds,
                decorations.pathreplacing, calc, matrix, trees}
\usepackage{pgfplots}
\pgfplotsset{compat=1.18}

% --- Tables ---
\usepackage{booktabs}
\usepackage{multirow}
\usepackage{array}
\usepackage{tabularx}
\usepackage{makecell}

% --- Typography & Layout ---
\usepackage{enumitem}
\usepackage{parskip}
\usepackage{fancyhdr}
\usepackage{titlesec}
\usepackage{titletoc}
\usepackage{hyperref}
\usepackage{graphicx}
\usepackage{float}
\usepackage{caption}

% --- Color Palette ---
\definecolor{defblue}{HTML}{1A5276}      % Definition header
\definecolor{defbluebg}{HTML}{D6EAF8}    % Definition background
\definecolor{thmgreen}{HTML}{1E8449}     % Theorem header
\definecolor{thmbg}{HTML}{D5F5E3}        % Theorem background
\definecolor{exorange}{HTML}{B7770D}     % Example header
\definecolor{exbg}{HTML}{FEF9E7}         % Example background
\definecolor{warnred}{HTML}{922B21}      % Warning/Key Insight header
\definecolor{warnbg}{HTML}{FADBD8}       % Warning background
\definecolor{noticegray}{HTML}{2C3E50}   % Note header
\definecolor{noticebg}{HTML}{EBF5FB}     % Note background
\definecolor{taketawayyellow}{HTML}{7D6608} % Key Takeaway header
\definecolor{takeawaybg}{HTML}{FDFEFE}   % Takeaway background
\definecolor{accentpurple}{HTML}{6C3483} % Accent / proof
\definecolor{lightgray}{HTML}{F2F3F4}    % shading

% ============================================================
%  TCOLORBOX ENVIRONMENTS
% ============================================================

% -- Definition (Blue) --
\newtcolorbox{defbox}[2][]{
  enhanced, breakable,
  colback=defbluebg, colframe=defblue, coltitle=white,
  fonttitle=\bfseries\sffamily,
  title={Definition: #2},
  attach boxed title to top left={yshift=-2mm, xshift=4mm},
  boxed title style={colback=defblue, rounded corners},
  arc=4pt, boxrule=1.2pt,
  left=6pt, right=6pt, top=8pt, bottom=6pt,
  #1
}

% -- Theorem / Proposition (Green) --
\newtcolorbox{thmbox}[2][]{
  enhanced, breakable,
  colback=thmbg, colframe=thmgreen, coltitle=white,
  fonttitle=\bfseries\sffamily,
  title={Theorem: #2},
  attach boxed title to top left={yshift=-2mm, xshift=4mm},
  boxed title style={colback=thmgreen, rounded corners},
  arc=4pt, boxrule=1.2pt,
  left=6pt, right=6pt, top=8pt, bottom=6pt,
  #1
}

% -- Proposition (Green variant) --
\newtcolorbox{propbox}[2][]{
  enhanced, breakable,
  colback=thmbg, colframe=thmgreen!70!black, coltitle=white,
  fonttitle=\bfseries\sffamily,
  title={Proposition: #2},
  attach boxed title to top left={yshift=-2mm, xshift=4mm},
  boxed title style={colback=thmgreen!70!black, rounded corners},
  arc=4pt, boxrule=1.2pt,
  left=6pt, right=6pt, top=8pt, bottom=6pt,
  #1
}

% -- Example (Orange/Yellow) --
\newtcolorbox{exbox}[2][]{
  enhanced, breakable,
  colback=exbg, colframe=exorange, coltitle=white,
  fonttitle=\bfseries\sffamily,
  title={Example: #2},
  attach boxed title to top left={yshift=-2mm, xshift=4mm},
  boxed title style={colback=exorange, rounded corners},
  arc=4pt, boxrule=1.2pt,
  left=6pt, right=6pt, top=8pt, bottom=6pt,
  #1
}

% -- Key Insight / Warning (Red) --
\newtcolorbox{warnbox}[2][]{
  enhanced, breakable,
  colback=warnbg, colframe=warnred, coltitle=white,
  fonttitle=\bfseries\sffamily,
  title={Key Insight: #2},
  attach boxed title to top left={yshift=-2mm, xshift=4mm},
  boxed title style={colback=warnred, rounded corners},
  arc=4pt, boxrule=1.2pt,
  left=6pt, right=6pt, top=8pt, bottom=6pt,
  #1
}

% -- Note (Gray-Blue) --
\newtcolorbox{notebox}[1][]{
  enhanced, breakable,
  colback=noticebg, colframe=noticegray!60, coltitle=white,
  fonttitle=\bfseries\sffamily,
  title={Note},
  attach boxed title to top left={yshift=-2mm, xshift=4mm},
  boxed title style={colback=noticegray!60, rounded corners},
  arc=4pt, boxrule=0.8pt,
  left=6pt, right=6pt, top=8pt, bottom=6pt,
  #1
}

% -- Key Takeaway --
\newtcolorbox{takebox}[1][]{
  enhanced, breakable,
  colback=takeawaybg, colframe=taketawayyellow, coltitle=white,
  fonttitle=\bfseries\sffamily,
  title={Key Takeaway},
  attach boxed title to top left={yshift=-2mm, xshift=4mm},
  boxed title style={colback=taketawayyellow, rounded corners},
  arc=4pt, boxrule=1.2pt,
  left=6pt, right=6pt, top=8pt, bottom=6pt,
  #1
}

% -- Intuition --
\newtcolorbox{intubox}[1][]{
  enhanced, breakable,
  colback=accentpurple!8, colframe=accentpurple!60, coltitle=white,
  fonttitle=\bfseries\sffamily,
  title={Intuition},
  attach boxed title to top left={yshift=-2mm, xshift=4mm},
  boxed title style={colback=accentpurple!80, rounded corners},
  arc=4pt, boxrule=0.8pt,
  left=6pt, right=6pt, top=8pt, bottom=6pt,
  #1
}

% ============================================================
%  PAYOFF MATRIX HELPER
%  Usage: inside tabular, use \payoff{a}{b} for (a,b)
% ============================================================
\newcommand{\payoff}[2]{$(#1,\;#2)$}
\newcommand{\nash}[2]{\underline{$#1$},\underline{$#2$}}  % underline NE payoffs
\newcommand{\Nset}{\mathbb{N}}
\newcommand{\Rset}{\mathbb{R}}
\newcommand{\Sset}{\mathcal{S}}
\newcommand{\Ui}{u_i}

% ============================================================
%  SECTION FORMATTING
% ============================================================
\titleformat{\section}
  {\large\bfseries\sffamily\color{defblue}}
  {\thesection}{1em}{}[\titlerule]

\titleformat{\subsection}
  {\normalsize\bfseries\sffamily\color{thmgreen!80!black}}
  {\thesubsection}{1em}{}

\titleformat{\subsubsection}
  {\normalsize\itshape\sffamily\color{exorange!80!black}}
  {\thesubsubsection}{1em}{}

% ============================================================
%  HEADER / FOOTER
% ============================================================
\pagestyle{fancy}
\fancyhf{}
\fancyhead[L]{\sffamily\small\color{defblue} MTH636 -- Game Theory}
\fancyhead[R]{\sffamily\small\color{defblue} \leftmark}
\fancyfoot[C]{\sffamily\small\thepage}
\renewcommand{\headrulewidth}{0.4pt}

% ============================================================
%  HYPERREF SETUP
% ============================================================
\hypersetup{
  colorlinks=true,
  linkcolor=defblue,
  urlcolor=thmgreen,
  citecolor=accentpurple,
  bookmarksnumbered=true,
  pdftitle={MTH636 Game Theory Notes},
  pdfauthor={MTH636 Student Notes}
}

% ============================================================
%  MISC HELPERS
% ============================================================
\newcommand{\BR}{\text{BR}}        % Best Response
\newcommand{\NE}{\text{NE}}        % Nash Equilibrium
\newcommand{\SPNE}{\text{SPNE}}    % Subgame Perfect NE
\newcommand{\Dom}{\text{dom}}      % Dominance
\newcommand{\mixs}[1]{\sigma_{#1}} % Mixed strategy
\newcommand{\pures}[1]{s_{#1}}     % Pure strategy


% ============================================================
\begin{document}

% Title
\begin{center}
  {\LARGE\bfseries\sffamily\color{defblue} MTH636 -- Game Theory}\\[6pt]
  {\Large\sffamily Lecture 3: Dominated Strategies \& Iterated Elimination}\\[4pt]
\end{center}
\vspace{0.5cm}
\hrule height 1.5pt
\vspace{0.3cm}

% ============================================================
\section{Strictly Dominated Strategies}

\begin{defbox}{Strictly Dominated Strategy}
  A strategy $s_i$ of player $i$ is \textbf{strictly dominated} by another strategy $t_i$ of the same player if, for all possible strategy profiles of the other players $s_{-i} \in S_{-i}$, playing $t_i$ always yields a strictly higher payoff than playing $s_i$. Mathematically:
  \[
    u_i(s_i, s_{-i}) < u_i(t_i, s_{-i}) \quad \forall s_{-i} \in S_{-i}
  \]
\end{defbox}

\begin{warnbox}{Rationality Assumptions}
  \begin{itemize}[leftmargin=*]
    \item \textbf{Assumption 1:} A rational player will \emph{never} play a strictly dominated strategy.
    \item \textbf{Assumption 2:} All players in a game are rational.
  \end{itemize}
\end{warnbox}

Let $S_{-i} = \prod_{j \neq i} S_j$ denote the set of all possible combinations of strategies for all players other than $i$.

\subsection{Iterated Elimination of Strictly Dominated Strategies (IESDS)}

\begin{exbox}{IESDS Process}
  Consider the following game:
  \begin{center}
  \renewcommand{\arraystretch}{1.3}
  \begin{tabular}{c|ccc}
    $P_1 \setminus P_2$ & L & M & R \\
    \hline
    T & \payoff{1}{0} & \payoff{1}{2} & \payoff{0}{1} \\
    B & \payoff{0}{3} & \payoff{0}{1} & \payoff{2}{0} \\
  \end{tabular}
  \end{center}

  \textbf{Step 1:} Look at Player 2's strategies (L, M, R).
  \begin{itemize}
    \item If P1 plays T, P2 gets 0 (L), 2 (M), 1 (R).
    \item If P1 plays B, P2 gets 3 (L), 1 (M), 0 (R).
  \end{itemize}
  Notice that for Player 2, strategy \textbf{M strictly dominates R} because $2 > 1$ and $1 > 0$. Since Player 2 is rational, they will never play R. We eliminate column R:
  
  \begin{center}
  \renewcommand{\arraystretch}{1.3}
  \begin{tabular}{c|cc}
    $P_1 \setminus P_2$ & L & M \\
    \hline
    T & \payoff{1}{0} & \payoff{1}{2} \\
    B & \payoff{0}{3} & \payoff{0}{1} \\
  \end{tabular}
  \end{center}

  \textbf{Step 2:} Now look at Player 1's strategies (T, B) in the reduced game.
  \begin{itemize}
    \item T yields payoffs (1, 1).
    \item B yields payoffs (0, 0).
  \end{itemize}
  For Player 1, strategy \textbf{T strictly dominates B} ($1 > 0$ and $1 > 0$). Player 1 knows Player 2 is rational and won't play R. Thus, we eliminate row B:

  \begin{center}
  \renewcommand{\arraystretch}{1.3}
  \begin{tabular}{c|cc}
    $P_1 \setminus P_2$ & L & M \\
    \hline
    T & \payoff{1}{0} & \payoff{1}{2} \\
  \end{tabular}
  \end{center}

  \textbf{Step 3:} Back to Player 2. In this further reduced game:
  \begin{itemize}
    \item L yields 0.
    \item M yields 2.
  \end{itemize}
  Strategy \textbf{M strictly dominates L} ($2 > 0$). Eliminate column L.

  \begin{center}
  \renewcommand{\arraystretch}{1.3}
  \begin{tabular}{c|c}
    $P_1 \setminus P_2$ & M \\
    \hline
    T & \payoff{1}{2} \\
  \end{tabular}
  \end{center}
  
  \textbf{Result:} The unique prediction is that Player 1 plays T and Player 2 plays M. This relies on Player 1 knowing that Player 2 is rational.
\end{exbox}

% ============================================================
\section{Common Knowledge and Complete Information}

\begin{defbox}{Common Knowledge}
  A fact is \textbf{Common Knowledge} among the players in a game if for any finite chain of players $i_1, i_2, \dots, i_k$, the following holds:
  Player $i_1$ knows that Player $i_2$ knows that $\dots$ Player $i_k$ knows the fact.
\end{defbox}

\begin{warnbox}{Assumption 3}
  \textbf{Rationality of the players is common knowledge.}
\end{warnbox}

\begin{defbox}{Game of Complete Information}
  A game is called a game of \textbf{complete information} if the structure of the game (the players, their strategies, and their payoff functions) is common knowledge among all players.
\end{defbox}

\begin{thmbox}{Order Independence of IESDS}
  The order of elimination of \emph{strictly} dominated strategies does not matter, in the sense that the final game obtained after removing all strictly dominated strategies does not depend on the order of elimination.
\end{thmbox}

% ============================================================
\section{Weakly Dominated Strategies}

\begin{defbox}{Weakly Dominated Strategy}
  A strategy $s_i$ of player $i$ is \textbf{weakly dominated} by another strategy $t_i$ of the same player if:
  \begin{enumerate}
      \item $u_i(s_i, s_{-i}) \leq u_i(t_i, s_{-i}) \quad \forall s_{-i} \in S_{-i}$
      \item $\exists t_{-i} \in S_{-i}$ such that $u_i(s_i, t_{-i}) < u_i(t_i, t_{-i})$
  \end{enumerate}
  (i.e., $t_i$ is always at least as good as $s_i$, and is strictly better against at least one strategy profile of the opponents).
\end{defbox}

\subsection{Iterated Elimination of Weakly Dominated Strategies}

Unlike strict domination, the order of elimination \emph{does} matter for weakly dominated strategies. 

\begin{exbox}{Order Dependence Example}
  Consider the following game:
  \begin{center}
  \renewcommand{\arraystretch}{1.3}
  \begin{tabular}{c|ccc}
    $P_1 \setminus P_2$ & L & C & R \\
    \hline
    T & \payoff{1}{2} & \payoff{2}{3} & \payoff{0}{3} \\
    M & \payoff{2}{2} & \payoff{2}{1} & \payoff{3}{2} \\
    B & \payoff{2}{1} & \payoff{0}{0} & \payoff{1}{0} \\
  \end{tabular}
  \end{center}
  
  Let's observe different valid elimination sequences:

  \textbf{Sequence 1: Eliminate T, R, B, C}
  \begin{enumerate}
      \item For P1, \textbf{M weakly dominates T} (since $2 > 1$, $2 = 2$, and $3 > 0$). Eliminate T.
      \item For P2, against \{M, B\}, \textbf{L weakly dominates R} (payoffs 2,1 vs 2,0). Eliminate R.
      \item For P1, against \{L, C\}, \textbf{M strictly dominates B} (payoffs 2,2 vs 2,0). Eliminate B.
      \item For P2, against \{M\}, \textbf{L strictly dominates C} (payoffs 2 vs 1). Eliminate C.
      \item \textbf{Result:} (M, L).
  \end{enumerate}

  \textbf{Sequence 2: Eliminate B, L, C, T}
  \begin{enumerate}
      \item For P1, \textbf{M weakly dominates B} (since $2 = 2$, $2 > 0$, $3 > 1$). Eliminate B.
      \item For P2, against \{T, M\}, \textbf{R weakly dominates L} (payoffs 3,2 vs 2,2). Eliminate L.
      \item For P2, against \{T, M\}, \textbf{R weakly dominates C} (payoffs 3,2 vs 3,1). Eliminate C.
      \item For P1, against \{R\}, \textbf{M strictly dominates T} (payoffs 3 vs 0). Eliminate T.
      \item \textbf{Result:} (M, R).
  \end{enumerate}
\end{exbox}

\begin{takebox}
  While the Iterated Elimination of Strictly Dominated Strategies (IESDS) yields a unique result regardless of the order of elimination, the \textbf{Iterated Elimination of Weakly Dominated Strategies is order-dependent}. Different elimination paths can lead to entirely different predictions.
\end{takebox}

\end{document}